\section{Einleitung}
\subsection{Zweck des Dokuments}
Dieses Dokument spezifiziert vollständig, eindeutig und überprüfbar das Nothaltesystem eines autonom fahrenden Miniatur-LKWs,
 das im Rahmen der Projektarbeit im Modul Software Engineering an der DHBW Friedrichshafen entwickelt wurde.
Ziel des Dokuments ist es, alle funktionalen und nicht-funktionalen Anforderungen zu definieren, 
das Systemverhalten in unterschiedlichen Betriebszuständen klar darzustellen und die Interaktion zwischen Softwarekomponenten, 
Sensorsystemen und Aktorik präzise zu beschreiben.
Die Spezifikation dient als Grundlage für Implementierung, Test, Validierung sowie sicherheitsrelevante Bewertung des Systems.
\subsection{Geltungsbereich des Systems}
Das in diesem Dokument beschriebene Nothaltesystem umfasst ausschließlich sicherheitsrelevante Funktionen zur Erkennung gefährlicher 
Situationen und zur unmittelbaren Einleitung geeigneter Reaktionen.
Dazu gehören:
\begin{itemize}
	\item \ac{Hardstop}
	\item \ac{Softstop}
	\item \ac{E-Stop}
	\item Quittierung und Wiederanlauf (\ac{Resume})
	\item \ac{Kriechmodus}
	\item Warnsignale (\ac{LED}, Dachleuchte)
	\item Protokollierung sicherheitsrelevanter Ereignisse
\end{itemize}
Nicht Bestandteil des Systems sind:
\begin{itemize}
	\item allgemeine Fahrzeugsteuerung
	\item Navigationsalgorithmen
	\item Hindernisvermeidung
	\item Odometrie
	\item Fahrdynamikregelung
\end{itemize}
Das Nothaltesystem arbeitet ergänzend zum autonomen Fahrsystem, besitzt jedoch Vorrang bei sicherheitskritischen Situationen.




